% Bewertungspaket zum Gesamttest «Polynomfunktionen» — Quelle des PDFs (python3 scripts/build-lp-pdf.py).
% Ganze Punkte je Teilschritt; Werte mit python3 nachgerechnet (04.10.2026, Folgefehler-Fälle eingeschlossen; nach der Prüfung überarbeitet).
\documentclass[11pt]{article}
\usepackage{../lp-druck}
\renewcommand{\lpname}{Polynomfunktionen}
\renewcommand{\lpbereich}{Schwerpunktfach}
\renewcommand{\lpteilgebiet}{3.3}
\renewcommand{\lpfuss}{Bewertungspaket}
\newcolumntype{P}{>{\centering\arraybackslash}p{10mm}}
\newenvironment{raster}{\par\smallskip\noindent\tabularx{\linewidth}{@{}>{\raggedright\arraybackslash}p{38mm}P X@{}}\toprule
  \small\textsc{Teilschritt} & \small\textsc{P} & \small\textsc{Musterlösung}\\\midrule}{\bottomrule\endtabularx\par}
\newcommand{\fehler}[1]{\par\smallskip{\small\textcolor{lprot}{\bfseries Typische Fehler:} #1}}
\newcommand{\E}{\,{\footnotesize\bfseries\color{lpgruen}(E)}}
\newcommand{\A}{\,{\footnotesize\bfseries\color{lpblau}(A)}}
\begin{document}
\lpkopf{Bewertungspaket zum Gesamttest}{}

\begin{lpachtung}{Erst lösen, dann öffnen}
Dieses Paket enthält die Musterlösung. Wer es vor dem Lösen liest, misst am Ende nur, wie gut das Abschreiben gelingt.
\end{lpachtung}

\section*{1 · So gehst du vor}
\begin{enumerate}[leftmargin=*,itemsep=2pt]
\item \textbf{Gesamttest lösen} — vollständig, mit Rechenweg, ohne dieses Paket. Teil~A ohne, Teil~B mit Taschenrechner.
\item \textbf{Fotografieren oder scannen}, gut lesbar, eine Seite pro Bild. \textbf{Kein Name} auf den Bildern.
  Handyfotos tragen oft unsichtbar Standort, Zeit und Gerät mit (Metadaten). Lade darum lieber einen Scan oder ein
  Bildschirmfoto des Fotos hoch, oder entferne vorher die Standortangaben.
\item \textbf{Bewerten lassen.} Ob du dafür eine KI nutzt und welche, entscheidest du selbst und in eigener Verantwortung —
  je nachdem, was dir aufgrund deines Alters und deiner persönlichen Situation erlaubt ist (Altersgrenzen und
  Nutzungsbedingungen des Anbieters, allenfalls Einverständnis deiner Eltern). Lade nur deine Lösung und dieses PDF hoch,
  keine persönlichen Daten, und füge den Auftrag aus Abschnitt~2 ein. Ohne KI kannst du dich mit Abschnitt~3 selbst bewerten.
\item \textbf{Rückmeldung prüfen.} Stimmt, was die KI aus deiner Handschrift gelesen hat? Eine KI kann sich irren —
  im Zweifel gilt das Raster in Abschnitt~3.
\item \textbf{Gezielt wiederholen} nach Abschnitt~4.
\end{enumerate}

\needspace{14\baselineskip}
\section*{2 · Auftrag an die KI}
\begin{tcolorbox}[colback=lplinie!25,colframe=lplinie,boxrule=0.5pt,arc=1mm,fontupper=\small]
Du bist Korrektorin oder Korrektor für Mathematik an einer Schweizer Berufsmaturitätsschule. Im Anhang findest du (1)~das Bewertungspaket zum Gesamttest «Polynomfunktionen» mit Musterlösung und Punkteraster und (2)~Fotos meiner handschriftlichen Lösung.\par\medskip
Geh so vor:\par
1. Schreib zu jeder Aufgabe G1 bis G8 zuerst ab, was du in meiner Lösung liest — Rechenweg und Ergebnis. Was du nicht sicher lesen kannst, markierst du mit [unsicher]. Nicht raten.\par
2. Vergib die Punkte genau nach dem Raster im Bewertungspaket (Abschnitt 3), ganze Punkte je Zeile. Ziehe einen Fehler nur einmal ab: Rechne ich mit einem falschen Zwischenergebnis richtig weiter, gibt es die Folgepunkte. Ausnahme: Zeilen mit \textbf{(E)} gibt es nur für das richtige Ergebnis, nie als Folgepunkt. Die Zeilen «Typische Fehler» sagen, welche Rasterzeile bei diesem Fehler 0 Punkte gibt — sie sind kein zusätzlicher Abzug.\par
3. Andere richtige Lösungswege zählen voll. Gleichwertige Schreibweisen zählen voll: Dezimalpunkt oder Dezimalkomma, Bruch oder Dezimalzahl ($-\tfrac12 = -0.5$), $y = \ldots$ oder $f(x) = \ldots$, Faktoren in anderer Reihenfolge, $(x-2)(x-2)$ statt $(x-2)^2$ — ausser wo das Raster eine bestimmte Form verlangt. Drei Stufen: Zeilen mit \textbf{(A)} verlangen nur Ablesen, diese Punkte gibt es auch ohne Rechenweg. Zeilen mit \textbf{(E)} sind Ergebnispunkte: Das Ergebnis muss stimmen \emph{und} aus einem erkennbaren Weg hervorgehen, ausser die Zeile sagt ausdrücklich etwas anderes. Alle übrigen Zeilen bewerten den Weg selbst.\par
4. Begründe jeden Abzug in einem Satz und nenne die Fehlerart (zum Beispiel «Vorzeichen in der Klammer» oder «Quadrat beim Einsetzen vergessen»). Schreib die Musterlösung nicht ab — zeig mir, wo mein Fehler liegt.\par
5. Gib zum Schluss eine Tabelle «Aufgabe | Punkte | Begründung», die Gesamtpunktzahl (von 25) und — nach der Selbsteinschätzung in Abschnitt 4 — welches Kapitel des Leitprogramms ich wiederholen soll. Keine Note.\par
6. Fehlt ein Foto oder ist eine Aufgabe unlesbar, sag es, statt sie zu bewerten.
\end{tcolorbox}

\needspace{16\baselineskip}
\section*{3 · Musterlösung und Punkteraster}
\textbf{Allgemein:} Ganze Punkte je Zeile. Ein Fehler wird nur einmal abgezogen (Folgefehler). Gleichwertige Wege und
Schreibweisen zählen voll. In der Spalte P:
\textbf{(A)}~= nur ablesen, zählt auch ohne Rechenweg · \textbf{(E)}~= Ergebnispunkt, Ergebnis richtig
\emph{und} Weg erkennbar (ausser die Zeile sagt ausdrücklich etwas anderes) · ohne Zeichen = die Zeile bewertet den Weg.
Teil~A ist ohne Taschenrechner gestellt, Teil~B mit.
«Typische Fehler» nennt, welche Zeile 0 Punkte gibt, und die Punkte, die bleiben.
\textbf{Teilrichtige Zeilen:} Verlangt eine Zeile mehrere Angaben und stimmt nur ein Teil, gibt die Zeile 0 Punkte —
ausser die Zeile nennt ausdrücklich eine Aufteilung.

\begin{lpaufgabe}{G1}{$f(x) = -0.5\,(x+3)\,(x-2)^{3}$}{4 P}
\begin{raster}
(a) Grad und Leitkoeffizient & 1\A & Grad $1 + 3 = 4$, Leitkoeffizient $-0.5$. Beides nötig.\\
(b) Nullstellen mit Vielfachheit & 1\A & $-3$ einfach, $2$ dreifach. Beides nötig.\\
(b) Verhalten & 1 & Bei $-3$ schneidet der Graph die $x$-Achse; bei $2$ schneidet er sie mit einer Terrasse (flach, Vorzeichenwechsel).\\
(c) $f(0)$ & 1\E & $f(0) = -0.5 \cdot 3 \cdot (-2)^{3} = -0.5 \cdot 3 \cdot (-8) = 12$\\
\end{raster}
\fehler{Grad 2 («zwei Klammern») oder Grad 3: Zeile (a) 0 $\to$ 3 von 4\,P.
Nullstellen $3$ und $-2$ (Vorzeichen in der Klammer): Zeile (b) erste 0; passt das Verhalten zur eigenen Vielfachheit, zählt die Verhaltens-Zeile $\to$ 3 von 4\,P.
«bei $2$ berührt er»: Verhaltens-Zeile 0 $\to$ 3 von 4\,P.
$f(0) = 3$ (Exponent vergessen: $-0.5 \cdot 3 \cdot (-2)$): Zeile (c) 0 $\to$ 3 von 4\,P.
$f(0) = -12$ ($(-2)^3$ als $+8$ gerechnet): Zeile (c) 0 $\to$ 3 von 4\,P.}
\end{lpaufgabe}

\begin{lpaufgabe}{G2}{Graph $\to$ Gleichung}{4 P}
\begin{raster}
Nullstellen abgelesen & 1\A & $-2$ und $2$\\
Vielfachheit begründet & 1 & Bei $2$ berührt der Graph die $x$-Achse (er bleibt auf derselben Seite): doppelte Nullstelle. Bei $-2$ schneidet er: einfach. Ansatz $f(x) = a\,(x+2)\,(x-2)^{2}$.\\
$a$ aus $(0 \mid 4)$ & 1 & $f(0) = a \cdot 2 \cdot (-2)^{2} = 8a = 4$\\
Gleichung & 1\E & $a = 0.5$, also $f(x) = 0.5\,(x+2)\,(x-2)^{2}$\\
\end{raster}
\fehler{Ansatz $a\,(x+2)^{2}(x-2)$ (Vielfachheit an der falschen Stelle): Begründungs-Zeile 0; damit $f(0) = -8a = 4$, $a = -0.5$ — die $a$-Zeile zählt als Folgefehler, die Gleichung (E) nicht $\to$ 2 von 4\,P.
Quadrat beim Einsetzen vergessen ($a \cdot 2 \cdot (-2) = -4a = 4$, $a = -1$): $a$-Zeile 0, Gleichung 0 $\to$ 2 von 4\,P. Hinweis: Rechts geht der Graph nach oben, also muss $a$ positiv sein.
Nullstellen mit falschem Vorzeichen in den Klammern ($a\,(x-2)(x+2)^2$ als Ansatz für «$2$ doppelt»): Begründungs-Zeile 0; $f(0) = -8a = 4$, $a = -0.5$ zählt als Folgefehler, die Gleichung nicht $\to$ 2 von 4\,P.}
\end{lpaufgabe}

\begin{lpaufgabe}{G3}{$g(x) = -x^{4} + 5x^{2} - 4$}{3 P}
\begin{raster}
(a) Enden & 1 & Grad $4$ gerade, Leitkoeffizient $-1 < 0$: beide Enden zeigen nach unten. Aussage und Begründung nötig.\\
(b) Symmetrie & 1 & Exponenten $4$, $2$ und $0$ (das konstante Glied) — alle gerade: achsensymmetrisch zur $y$-Achse. Ein Nachweis mit $g(-x) = g(x)$ zählt ebenso.\\
(c) Höchstzahlen & 1 & höchstens $4$ Nullstellen und höchstens $3$ lokale Extremstellen. Beides nötig.\\
\end{raster}
\fehler{«beide Enden oben» (Vorzeichen übersehen) oder «links oben, rechts unten» (Parität verwechselt): Zeile (a) 0 $\to$ 2 von 3\,P.
«weder noch, wegen $-4$»: Zeile (b) 0 — das konstante Glied hat den Exponenten $0$, und der ist gerade $\to$ 2 von 3\,P.
«höchstens $4$ Extremstellen»: Zeile (c) 0 $\to$ 2 von 3\,P.}
\end{lpaufgabe}

\begin{lpaufgabe}{G4}{Zweimal berühren}{2 P}
\begin{raster}
Mindestgrad & 1 & Grad mindestens $4$: Berühren heisst gerade Vielfachheit, also mindestens doppelt — zwei doppelte Nullstellen geben $2 + 2 = 4$.\\
Vorzeichen & 1 & $a_n < 0$: Beide Enden zeigen nach unten, das gibt es nur bei geradem Grad mit negativem Leitkoeffizienten. Beispiel: $-(x+1)^2(x-3)^2$.\\
\end{raster}
\fehler{«Grad 2», weil es zwei Nullstellen sind: erste Zeile 0 $\to$ höchstens 1 von 2\,P.
Nur «Grad 4» oder nur «negativ» ohne Begründung: die betroffene Zeile 0.}
\end{lpaufgabe}

\begin{lpaufgabe}{G5}{$h(x) = x^{3} - 4x^{2} + x + 6$}{4 P}
\begin{raster}
Nullstelle geraten mit Probe & 1 & Teiler von $6$ probieren, z.\,B. $h(-1) = -1 - 4 - 1 + 6 = 0$. Auch $h(2) = 0$ oder $h(3) = 0$ zählt — jede richtige Probe.\\
Linearfaktor abgespalten & 1 & Ansatz $(x + 1)(x^2 + px + q) = x^3 + (p + 1)x^2 + (q + p)x + q$: $p = -5$, $q = 6$, Kontrolle $q + p = 1$ (stimmt). Oder Polynomdivision $(x^{3} - 4x^{2} + x + 6) : (x + 1) = x^{2} - 5x + 6$. Beide Wege zählen voll (andere Startnullstelle: entsprechender Quotient, z.\,B. mit $x - 2$: $x^2 - 2x - 3$).\\
Quotient gelöst & 1 & $x^{2} - 5x + 6 = (x-2)(x-3)$ oder Lösungsformel\\
Nullstellen und Darstellung & 1\E & $-1$, $2$, $3$; $h(x) = (x+1)(x-2)(x-3)$\\
\end{raster}
\fehler{Nur mit dem Rechner (Gleichungslöser) die Nullstellen ohne Weg: die drei Weg-Zeilen 0, Ergebnis zählt nicht als (E) $\to$ 0 von 4\,P. Mit Rechner \emph{kontrolliert} ist erlaubt.
Probe falsch gerechnet ($h(1) = 0$ behauptet; tatsächlich $h(1) = 4$) und dann mit $x - 1$ weitergerechnet: Probe-Zeile 0; die Abspalt-Zeile 0, weil die Kontrolle (bzw. der Rest $4$ der Division) den Fehler zeigt und übergangen wurde; wird der so entstandene Quotient (z.\,B. $x^2 - 3x - 2$) folgerichtig gelöst, zählt die Quotienten-Zeile als Folgefehler; Ergebnis 0 $\to$ 1 von 4\,P.
Quotient falsch (z.\,B. $x^2 - 3x + 6$): Abspalt-Zeile 0; der Quotient wird als Folgefehler bewertet, die Ergebnis-Zeile nicht $\to$ höchstens 2 von 4\,P.
Nullstellen richtig, Darstellung fehlt oder Vorzeichen in den Klammern falsch: letzte Zeile 0 $\to$ 3 von 4\,P.}
\end{lpaufgabe}

\begin{lpaufgabe}{G6}{Am Graphen: $k(x) = -x^{3} + 3x^{2} + 1$}{3 P}
\begin{raster}
(a) Hoch- und Tiefpunkt & 1\A & $H(2 \mid 5)$, $T(0 \mid 1)$. Beide nötig.\\
(b) absolutes Maximum & 1 & $5$ im Hochpunkt $x = 2$: Der linke Randwert $k(-0.5) = 1.875$ ist kleiner.\\
(b) absolutes Minimum & 1\E & Randwert $k(3.5) = -42.875 + 36.75 + 1 = -5.125$; kleiner als $T$ (1) und als $k(-0.5)$: absolutes Minimum $-5.125$ am rechten Rand $x = 3.5$.\\
\end{raster}
\fehler{«Das absolute Minimum ist der Tiefpunkt $T(0 \mid 1)$»: Ränder nicht geprüft — letzte Zeile 0 $\to$ 2 von 3\,P.
$H$ und $T$ vertauscht: Zeile (a) 0; (b) als Folgefehler bewerten, die (E)-Zeile nur bei richtigem Wert $\to$ höchstens 2 von 3\,P.
Ein am Graphen abgelesener Randwert «etwa $-5$» ohne Rechnung: letzte Zeile 0 (Ergebnispunkt verlangt den Weg) $\to$ 2 von 3\,P.}
\end{lpaufgabe}

\begin{lpaufgabe}{G7}{$p(x) = 0.5x^{2} - 3x + 2$}{2 P}
\begin{raster}
$x_s$ & 1 & $x_s = -\dfrac{b}{2a} = -\dfrac{-3}{2 \cdot 0.5} = 3$ (oder über die quadratische Ergänzung)\\
$y_s$ und Art & 1\E & $p(3) = 4.5 - 9 + 2 = -2.5$; $a = 0.5 > 0$: Tiefpunkt $T(3 \mid -2.5)$. Beides nötig.\\
\end{raster}
\fehler{$x_s = -3$ (Minus vergessen): erste Zeile 0; der Punkt ist dann falsch, die (E)-Zeile 0 $\to$ 0 von 2\,P.
«Hochpunkt»: zweite Zeile 0 $\to$ 1 von 2\,P.
$p(3) = -11.5$ ($-3 \cdot 3$ als $+9$ … o.\,ä.) oder ein anderer Rechenfehler beim Einsetzen: zweite Zeile 0 $\to$ 1 von 2\,P.}
\end{lpaufgabe}

\begin{lpaufgabe}{G8}{Flugbahn $h(t) = -t^{3} + 6t^{2}$}{3 P}
\begin{raster}
(a) Nullstellen durch Ausklammern & 1 & $-t^{3} + 6t^{2} = -t^{2}\,(t - 6) = 0$: $t = 0$ doppelt, $t = 6$ einfach. Ausklammern, beide Nullstellen und die Vielfachheit von $0$ nötig.\\
(b) Wertetabelle & 1 & mindestens um das Maximum: $h(3) = 27$, $h(4) = 32$, $h(5) = 25$ (vollständig: $0, 5, 16, 27, 32, 25, 0$)\\
(b) Höchste Stelle & 1\E & In der Tabelle am höchsten bei $t = 4$ s mit $32$ m.\\
\end{raster}
\fehler{Nur $t = 6$ (Faktor $t^2$ übersehen) oder durch $t^2$ geteilt: Zeile (a) 0 $\to$ 2 von 3\,P.
$0$ als «einfach» angegeben: Zeile (a) 0 $\to$ 2 von 3\,P.
Tabellenwerte falsch, aber das Maximum folgerichtig aus der eigenen Tabelle: Tabellen-Zeile 0, die (E)-Zeile ebenfalls 0 $\to$ 1 von 3\,P.
Ohne Einheiten: kein Abzug.}
\end{lpaufgabe}

\needspace{14\baselineskip}
\section*{4 · Selbsteinschätzung}
\begin{tabularx}{\linewidth}{@{}l X@{}}\toprule
22 – 25 P & Die geprüften Teile sitzen. Wo du Punkte verloren hast: das Kapitel dieser Aufgabe nochmals (Zuordnung unten).\\
17 – 21 P & Den schwächsten Teil nochmals: Simulation und Übungen des Kapitels, in dem du die meisten Punkte verloren hast.\\
11 – 16 P & Zurück zu den Kapiteln aller Aufgaben, in denen du Punkte verloren hast.\\
0 – 10 P & Zurück zu Kapitel 1 und von dort der Reihe nach weiter.\\\midrule
\multicolumn{2}{@{}p{\linewidth}@{}}{\small\textbf{Aufgabe $\to$ Kapitel:} G1 $\to$ Kapitel 1 und 2; G2 $\to$ Kapitel 2; G3, G4 $\to$ Kapitel 3 (G4 auch 2); G5 $\to$ Kapitel 4; G6, G7, G8 $\to$ Kapitel 5}\\\bottomrule
\end{tabularx}
\end{document}
